% BHU PYQ -- Professional Exam Template
\documentclass[11pt,a4paper]{article}

\usepackage[a4paper,top=2.5cm,bottom=2.5cm,left=2.8cm,right=2.8cm,headheight=14pt]{geometry}
\usepackage{amsmath,amssymb}
\usepackage[T1]{fontenc}
\usepackage[utf8]{inputenc}
\usepackage{lmodern}
\usepackage{microtype}
\usepackage{fancyhdr}
\usepackage{enumitem}
\usepackage{listings}
\usepackage{hyperref}
\usepackage{lastpage}
\usepackage{parskip}

\hypersetup{colorlinks=true,urlcolor=black,linkcolor=black,
  pdftitle={BCS-403: Fundamentals Of Computing -- BHU PYQ}}

\pagestyle{fancy}
\fancyhf{}
\renewcommand{\headrulewidth}{0.4pt}
\renewcommand{\footrulewidth}{0.4pt}
\fancyhead[L]{\small\textit{Banaras Hindu University}}
\fancyhead[R]{\small\textit{BCS-403: Fundamentals Of Computing}}
\fancyfoot[C]{\small Page \thepage\ of \pageref{LastPage}}

\lstset{
  basicstyle=\small\ttfamily,
  frame=single,
  breaklines=true,
  columns=flexible,
  keepspaces=true
}

\newlist{parts}{enumerate}{2}
\setlist[parts,1]{label=(\alph*),leftmargin=*,itemsep=4pt,topsep=3pt}
\setlist[parts,2]{label=(\roman*),leftmargin=*,itemsep=2pt,topsep=2pt}
\newcommand{\pts}[1]{\hfill\textit{[#1]}}

\begin{document}
\begin{center}
  {\large\bfseries BANARAS HINDU UNIVERSITY}\\[2pt]
  {\normalsize B.Sc. (Hons.) Semester IV Examination, 2013-14}\\[6pt]
  \rule{0.8\linewidth}{0.5pt}\\[6pt]
  {\large\bfseries Paper: BCS-403 --- Fundamentals Of Computing}\\[4pt]
  {\normalsize Time: Three hours \hfill Maximum Marks: 70}
\end{center}

\rule{\linewidth}{0.4pt}

\smallskip
\noindent\textit{Instructions: Answer any five questions, including question no. 1, which is compulsory.}

\bigskip

\medskip\noindent\textbf{Question 1.} Write the full name and function of the following terms: \pts{$2\times7=14$}
\begin{parts}
  \item UNIVAC I
  \item ULSI
  \item CD-ROM
  \item ARPANET
  \item FTP
  \item URL
  \item HTTP
\end{parts}

\medskip\noindent\textbf{Question 2.}
\begin{parts}
  \item What is computer and what are its characteristics? Explain each. \pts{7}
  \item Differentiate among computer program, software and hardware. Explain the system software and application software with appropriate examples. \pts{2+5}
\end{parts}

\medskip\noindent\textbf{Question 3.}
\begin{parts}
  \item Divide $(1010110)_2$ by $(110)_2$. \pts{3}
  \item Multiply $(10101)_2$ with $(101)_2$. \pts{2}
  \item Add these binary numbers: $11111, 11111, 10111, 1011, 10111$. \pts{3}
  \item Convert the binary number $11011101.0111_2$ to decimal number system. \pts{3}
  \item Convert the decimal number $352_{10}$ to binary number system. \pts{3}
\end{parts}

\medskip\noindent\textbf{Question 4.}
\begin{parts}
  \item What do you mean by software development steps? Explain each step. \pts{7}
  \item What is programming language? Describe the different types of programming language with merits and demerits of each. \pts{7}
\end{parts}

\medskip\noindent\textbf{Question 5.}
\begin{parts}
  \item What is Internet? How is it different from the World Wide Web? Write the various applications of Internet. \pts{2+2+3}
  \item What are the steps involved in solving problems using computers? Explain each. \pts{7}
\end{parts}

\medskip\noindent\textbf{Question 6.}
\begin{parts}
  \item Write the characteristics of an algorithm. Draw a flowchart to check if a given year is a leap year or not. \pts{3+4}
  \item What is pseudocode? Write a pseudocode to find the sum of digits of a given number. \pts{2+5}
\end{parts}

\medskip\noindent\textbf{Question 7.} Write short notes on the following:
\begin{parts}
  \item Evolution of Computers \pts{4}
  \item Structured Programming and Logic/Control Structures \pts{5}
  \item Flowchart with its merits and demerits. \pts{5}
\end{parts}

\vfill
\rule{\linewidth}{0.4pt}
\begin{center}\small
\href{https://bhu-exam.vercel.app/}{Click here to visit our website}\\[4pt]\href{https://me-aryan.vercel.app/}{Click here to visit our contributor}\\[4pt]\href{https://quashan-me.vercel.app/}{Click here to visit our contributor}
\end{center}
\end{document}